En el análisis y diseño de algoritmos, comparar experimentalmente las cotas de complejidad teórica con el
 rendimiento empírico es fundamental para comprender el impacto real de las constantes ocultas y la sobrecarga 
 computacional~\cite{cormen2009,hales2017replicability}. Este informe aborda el estudio comparativo de dos problemas
  computacionales clásicos: primero el ordenamiento de arreglos unidimensionales empleando
   \textit{Merge Sort}, \textit{Quick Sort} (optimizado con partición de Hoare~\cite{educative_hoare_lomuto} 
   y eliminación por recursión de cola), \textit{Patience Sort}~\cite{taocv3,stablesort2019lis} y el 
   algoritmo \textit{Sort} default. 
   Además la multiplicación de matrices cuadradas comparando el método clásico $\mathcal{O}(N^3)$ frente al 
   método de partición en bloques de Strassen $\mathcal{O}(N^{2.807})$~\cite{cormen2009,bytequest2023strassen}. 
   El objetivo central es contrastar las curvas empíricas de ejecución frente a sus órdenes asintóticos ideales 
   bajo diversas familias de entradas y orden de magnitud de datos.